\section{The retained densities and their locality norms}
\label{sec:setup}

The block transformation will retain a coupling, finitely many polynomial
coefficient lists, a regular local remainder, and a separate contribution
from configurations with large nearest-neighbour differences.  This section
defines these objects.  In particular, all later comparisons concern
complete densities, even though some terms in their representation may
have either sign.

We use the scalar block construction and the local representation
conventions of~\cite{OpenAI-O4}.  The definitions are given here for
spins on \(S^2\).  Their preservation under integration is the
\(O(3)\) assertion proved in Section~\ref{sec:rg}.  Neither that assertion
nor the continuum construction is an imported \(O(4)\) theorem.

\subsection{Lattices, reference scales, and the observation}
\label{setup:scales}

A lattice at one layer is a square grid with its current orthonormal
coordinate frame and unit nearest-neighbour spacing.  A finite lattice
is its quotient by a rank-two lattice of translations; the corresponding
infinite-grid formula will be called the \emph{bulk formula}.  At each step the period lattice is a sublattice of the lattice of block
translations, so the observation descends to the quotient.  The
shortest nonzero period, measured in the current lattice units, is denoted
by \(\mathcal L_j\).  We use rectangular periods and the bounded-aspect
oblique periods specified in Section~\ref{sec:free}.  Local distances may
be measured in the maximum or \(\ell^1\) norm; we specify \(\ell^1\) for
the coefficient norm below and otherwise use the maximum norm.

Choose a dyadic integer \(L\), eventually sufficiently large, and a
terminal reference coupling \(H\).  For integers \(j\ge0\), set
\begin{equation}
 \begin{aligned}
 \gamma&=\frac{\log L}{2\pi},&
 H_j&=H+\gamma j,&
 \mathfrak m_j&=\left\lceil(\log H_j)^2\right\rceil,\\
 p_j&=(\log H_j)^{P_0},&
 t_j&=H_j^{-1/2}p_j,&
 T_j&=H_j^{-49/100}.
 \end{aligned}
 \label{setup:scale-definitions}
\end{equation}
These are the scales of \Rref{rg:scales}, with the drift normalization
appropriate to \(S^2\).  A run of depth \(N\) visits the layers
\(N,N-1,\ldots,0\).  In a step from \(j\) to \(j-1\), the precise
coupling \(b\) lies in \([H_j/2,2H_j]\), and
\(g=b^{-1/2}\).  We abbreviate \(M=\mathfrak m_j\),
\(p=p_j\), and \(t=t_j\) within a step.  Thus \(t\asymp gp\).
The thresholds \(H_j,p_j,t_j,T_j\) remain fixed when two precise
couplings are compared.

A \emph{free history} records the scalar observations already performed,
including their coordinate frames.  Its length is \(N-j\) on layer
\(j\).  For regular observations the history can be identified with its
length; after an inclined first observation the record, and not only
its length, matters.  We write \(h\) for this history.  The allowed
inclinations and their common analytic estimates are established in
Section~\ref{sec:free}.

Constants written \(C\) when choosing \(L\) are uniform for all
sufficiently large \(L\).  Constants \(C_L,c_L>0\) may depend on the
subsequently fixed \(L\); after the stated choices they are independent
of the layer, free history, period, and varying precise coupling.  All requested derivative orders, support
exponents, and polynomial powers in the geometric estimates are fixed
before \(P_0\) and \(H\); we choose \(P_0\) sufficiently large and then
\(H\) sufficiently large.  Fixed powers of \(M\) and \(p\) therefore
remain powers of \(\log H_j\).  Admitted periods satisfy \(\mathcal L_j\ge C_{\rm adm}\mathfrak m_j\),
with a fixed constant large enough for the prescribed local neighbourhoods.
The minimum terminal period \(m_{\min}(L,H)\) is enlarged accordingly.  Since \(\mathfrak m_j/\mathfrak m_{j-1}\le2\) for
large \(H\), the same choice works at all layers of a run.

For a regular step, partition the fine grid into \(L\)-by-\(L\) blocks
\(Y\).  Let \(w_Y(x)\ge0\), \(x\in Y\), be the sampled and normalized
product of a fixed even smooth bump, centred at the block centre and
supported in the coordinate square of radius \(0.14L\).  It can be
chosen positive in the square of radius \(0.1L\).  Thus
\[
 \sum_{x\in Y}w_Y(x)=1,\qquad
 \sum_{x\in Y}w_Y(x)(x-x_Y)=0.
\]
Define the scalar averaging map and the spin mean by
\[
 (Q\phi)_Y=\sum_{x\in Y}w_Y(x)\phi_x,
 \qquad m_Y(q)=\sum_{x\in Y}w_Y(x)q_x.
\]
Independently for each block, let a tag \(\tau_Y\) select a site
\(x\in Y\) with probability \(w_Y(x)\), and put
\begin{equation}
 \widehat m_Y(q,\tau_Y)=
 \begin{cases}
 m_Y(q)/|m_Y(q)|,& |m_Y(q)|\ge\tfrac12,\\
 q_{\tau_Y},& |m_Y(q)|<\tfrac12.
 \end{cases}
 \label{setup:unit-mean}
\end{equation}
Tags may also be sampled where their values are unused.  With
\(\dd\omega\) denoting normalized area measure on \(S^2\), the exact
observation is the probability kernel
\begin{equation}
 \mathcal K_b(\dd V\mid q)
 =\E_{\tau}\prod_Y
 \frac{\exp\{-b|V_Y-\widehat m_Y(q,\tau_Y)|^2/2\}}
      {z_2(b)}\,\dd\omega(V_Y),
 \qquad
 z_2(b)=\int_{S^2}e^{-b|V-e|^2/2}\,\dd\omega(V).
 \label{setup:observation}
\end{equation}
Here \(e\) is any unit vector.  Rotation invariance makes the denominator
independent of \(e\); explicitly \(z_2(b)=(1-e^{-2b})/(2b)\) for
\(b>0\).  Thus integrating the density against
\(\mathcal K_b\) preserves its partition function exactly.  This is
\Rref{rg:observation} with observation precision one.  Its piecewise
mean is measurable everywhere; the proof will differentiate only
specified smooth extensions on graphs with small chords.

\subsection{The scalar quadratic carried by a history}
\label{setup:scalar-data}

The quadratic data in a history concern real scalar fields and are
independent of the spin dimension.  Choose one orientation for each nearest-neighbour bond, using the positive
coordinate directions.  Site and edge scalar products use counting measure,
and adjoints refer to those scalar products.  Start with the nearest-neighbour
precision \(K_0=d^*d\), where \((d\phi)_{xy}=\phi_y-\phi_x\).
If \(Q\) is the next averaging map, integrating the fine scalar field
in
\[
 \exp\left\{-\tfrac12\inner{\phi}{K_h\phi}
                   -\tfrac12\norm{V-Q\phi}_2^2\right\}
\]
gives the next precision
\begin{equation}
 K_{h+1}=\Id-Q(K_h+Q^*Q)^{-1}Q^*.
 \label{setup:scalar-schur}
\end{equation}
The inverse exists on each finite connected torus.  Inductively,
\(K_h\ge0\) has only constant null vectors, while \(Q\) preserves
constants; hence \(K_h+Q^*Q\) is positive definite.  The Schur complement
in \eqref{setup:scalar-schur} has only constant null vectors again, as is
seen by minimizing the displayed nonnegative quadratic form at fixed \(V\).  The formula defines the quadratic after extraction
of its scalar Gaussian normalization.  Bulk kernels are the common
translation-covariant kernels whose periodizations give these finite
operators.

For regular histories the scalar construction of \Rref{free:bounds}
and \Rref{lem:free-stack} supplies a fixed \(0<c_0<1\) and a real edge
operator \(\ell_h\) satisfying
\begin{equation}
 K_h=c_0d^*d+d^*\ell_h^*\ell_hd,
 \qquad \ell_0=\sqrt{1-c_0}\,\Id.
 \label{setup:edge-lift}
\end{equation}
Its kernels have exponential row and column moments, uniformly in the
regular history, with the constants and derivative bounds specified in
Section~\ref{sec:free}.  These are scalar statements: the operators act
componentwise on ambient vectors.  Section~\ref{sec:free} gives the
corresponding choices for the additional histories used here and proves
the stronger history comparison required later.  A retained density
always carries the particular history and edge operator chosen by these
common rules; it does not choose a new lift on each finite torus.

\subsection{The class of effective densities}
\label{subsec:setup-density}

At each layer, the exact integration retains a kinetic energy, finitely
many polynomial slots with summable coefficient lists, a differentiable
error, and a sum of
weights covering the edges with large spin differences.  We specify
these four parts and their norms here.  This is the density class of
\Rref{rg:class}, with the tangent coordinates and spatial symmetry
adapted to $S^2$.

Fix a finite layer torus $\Lambda_j$ and free history $h$.
Write $M=\mathfrak m_j$ and $t=t_j$.  Bulk formulas below are the
corresponding lifted formulas on the infinite lattice; the density
itself is defined on the finite torus.  All distances below are in the
current lattice units.  For an
edge $e=\langle x,y\rangle$, put
\[
 d_eq=q_y-q_x,\qquad r_e(q)=|d_eq|,
 \qquad D(q)=\{e:r_e(q)>t_j\}.
\]
The orientation chosen for $d_eq$ does not affect $r_e$.
Fix $0<d_0<c_0/16$ and choose $R_*$ sufficiently large compared with the
exponential localization length of the free edge operator $\ell_h$,
uniformly in the history.  Define
\begin{align}
 S_t(r)&=
 \begin{cases}
  c_0r^2/2,&0\le r\le t,\\
  d_0tr,&r>t,
 \end{cases}
 \label{eq:setup-clipped-energy}\\
 m_e(q)&=\mathbf 1\!\left\{
    r_f(q)\le t\exp\bigl(\operatorname{dist}(e,f)/R_*\bigr)
                     \text{ for every edge }f\right\},
 \label{eq:setup-kinetic-mask}\\
 \mathcal E_{h,t}(q)&=
       \sum_e S_t(r_e(q))
       +\frac12\sum_e m_e(q)| (\ell_hdq)_e|^2.
 \label{eq:setup-effective-energy}
\end{align}
These are the exact prescriptions of \Rref{rg:kinetic}; in particular,
$S_t$ is not replaced by a continuous interpolation at $r=t$.
Since $r_f\le2$, a mask $m_e$ queries only edges within distance
$O(\log(1/t))$ of $e$.  Its defining bound and the exponential row
moments give $| (\ell_hdq)_e|\le Ct$ whenever $m_e=1$.
The functions $S_t$ and $m_e$ can be
discontinuous.  The differentiable norms below apply to the retained
local functions on their stated open domains, rather than to these
indicators as functions of a moving output configuration.

\subsubsection{Supports, records, and graph masks}

A label includes its formula, its anchor when one is assigned, its
complete support, and a connected record covering that support.
Permissible anchors are chosen in the complete support and its record;
every anchored graph contains its anchor.  The complete support contains every spin argument and every edge or site
queried by a mask, an eligibility test, a test for the absence of another
object, or a window used to determine the formula.  A record consists of
cubes of radius $\mathfrak m_j$ together with the paths joining them.
It has an integer load $s\ge1$, chosen according to the conventions
following \Rref{rg:kinetic}, so that
\begin{equation}
 |\text{complete support}|\le C\mathfrak m_j^2s,
 \qquad
 \operatorname{length}(\text{record})\le C\mathfrak m_js.
 \label{eq:setup-record-size}
\end{equation}
The support cardinality here counts sites, including the endpoints of
queried edges.  Intersecting records can be joined with load at most a
fixed multiple of the sum of their loads.  The record retains all paths
used in a connected calculation, including multiple occurrences of a
path.

Upon projection to the next lattice and enlargement by the fixed
neighborhoods used in the integration, the load is chosen to satisfy
\begin{equation}
       \frac{4s}{L}\le s'\le D_*(1+s/L).
 \label{eq:setup-record-projection}
\end{equation}
The upper bound expresses the shortening of the joining paths under
blocking.  The lower bound is a convention: unused load is retained
when necessary.  The constant $D_*$ accommodates either the regular
blocking or the inclined first blocking.  This is the convention of
\Rref{rg:projection}.

On a torus, the record retains its lifted displacements before positions
are reduced modulo the period lattice.  Recall that $\mathcal L_j$ is
the shortest period length in the current lattice units.  A complete record is called
\emph{short} if
\begin{equation}
                         s<c_*\mathcal L_j/\mathfrak m_j.
 \label{eq:setup-short-record}
\end{equation}
The fixed constant $c_*>0$ is sufficiently small that a short record has
an injective lift and remains within the required bulk comparison after
the fixed geometric enlargements.  A short formula must agree with the
corresponding formula on the infinite lattice.  All other records are
called \emph{winding}.  This classification concerns the complete
calculation: a constant, or a function with small displayed support, is
still assigned a winding record if its construction used a long record
or a period-dependent discrepancy.  The lower bound in
\eqref{eq:setup-record-projection} retains the required load when an
already winding term is transferred.

For a connected graph $X$, define its mask and its open graph domain
by
\begin{equation}
 \begin{aligned}
 1_X(q)&=\mathbf 1\{r_e(q)\le t_j\text{ for every }e\in X\},\\
 \mathcal U_{X,j}&=\{q:r_e(q)<4t_j\text{ for every }e\in X\}.
 \end{aligned}
 \label{eq:setup-graph-domain}
\end{equation}
A graph may consist of a single site with no edges; in that case both
edge conditions in \eqref{eq:setup-graph-domain} are vacuous.
A graph used for a local function contains that function's spin
arguments and is included in its complete record.  A masked local
function is defined to be zero off its mask; no value of an undefined
extension is used there.

\subsubsection{Tangent polynomials and their coefficient norms}

For an anchor $o$, define the relative tangent coordinates and a
nearest-neighbor quadratic function by
\begin{equation}
 y_o(x)=q_x-(q_o\cdot q_x)q_o\in q_o^\perp,
 \qquad
 S_o(q)=\frac14\sum_{|e|_1=1}|y_o(o+e)|^2.
 \label{eq:setup-tangent-coordinates}
\end{equation}
Here $|\cdot|_1$ is the lattice $\ell^1$ distance in the current square
frame.  For $q_\circ\in S^2$, the spherical exponential is
$\operatorname{Exp}_{q_\circ}(u)=\cos|u|\,q_\circ+(\sin|u|/|u|)u$
for $u\in q_\circ^\perp$, with its continuous value at $u=0$.
Its restriction to $|u|<\pi$ is one-to-one onto
$S^2\setminus\{-q_\circ\}$; the inverse is denoted by
$\operatorname{Log}_{q_\circ}$.  The coordinate $y_o(x)$ is a globally
defined tangent projection, whereas $\operatorname{Log}_{q_o}q_x$ is
the principal geodesic coordinate on this domain.
For each required degree $n$, choose a fixed basis of
$((\R^2)^{\otimes n})^{O(2)}$.  A basis tensor is evaluated in
$q_o^\perp$ through any orthonormal identification of this plane with
$\R^2$.  Its $O(2)$ invariance makes the result independent of that
identification.  A degree-$n$ label is a coefficient $a$ multiplying
\[
 T\bigl(y_o(o+r_1),\ldots,y_o(o+r_n)\bigr),
 \qquad r_i\in\Z^2,
\]
together with its paths, mask, and complete record.  A label containing
a zero displacement is discarded, since $y_o(o)=0$.  The full $O(2)$
invariance forces every
odd-degree tensor to vanish.

The five potentially nonzero slots retain the seven-slot indexing of
\Rref{rg:slots}:
\begin{equation}
 (bP_4,\ I_2,\ bP_5,\ I_3,\ bP_6,\ I_4,\ b^{-1}J_2),
 \qquad P_5=I_3=0,
 \label{eq:setup-slots}
\end{equation}
and
$\mathbf P=(P_4,I_2,0,0,P_6,I_4,J_2)$ denotes the coefficient lists
without the displayed powers of $b$.  A subscript denotes homogeneous
degree in the relative tangent coordinates, rather than order in $g$.

For a position list, set
\begin{equation}
 u(r_1,\ldots,r_n)=1+\sum_{i=1}^n|r_i|_1,
 \qquad W_\sigma(u)=e^{\sigma u}(1+u)^2.
 \label{eq:setup-path-weight}
\end{equation}
The paths from an anchor to its arguments are counted with multiplicity.
Path rules are chosen equivariantly; when a symmetry exchanges tied
coordinate-order paths, the finite collection of choices is retained
together.  The complete record contains those choices.  The coefficient
path parameter in \eqref{eq:setup-path-weight} is the fixed convention
of \Rref{supp:canonical-path}.
For degree four or six, the norm of a coefficient list $A_n$ is
\begin{equation}
 \|A_n\|_\sigma
   =\sup_o\sum_{\text{labels at }o}|a|W_\sigma(u).
 \label{eq:setup-canonical-norm}
\end{equation}
The same fixed $\sigma>0$ is used at every layer; it is chosen sufficiently
small for the required free-kernel exponential moments, before choosing
$L$.

The quadratic norm also records the subtraction that removes the affine
Hessian.  Let $\mathcal C_4$ be the four quarter-turn rotations of the
current square frame and put
\begin{equation}
 M_{r,s}(q)=\frac14\sum_{\rho\in\mathcal C_4}
       y_o(o+\rho r)\cdot y_o(o+\rho s),
 \qquad
 Q_{r,s}(q)=M_{r,s}(q)-(r\cdot s)S_o(q).
 \label{eq:setup-quadratic-packet}
\end{equation}
The two factors of a quadratic monomial are symmetrized before taking
coefficient cancellations.  The spatial identity used for this
symmetrization is
\[
 \frac18\sum_{\rho\in\mathcal C_4}
 \bigl[(\rho r)_i(\rho s)_j+(\rho s)_i(\rho r)_j\bigr]
                 =\frac{r\cdot s}{2}\delta_{ij}.
\]
If $y_o(o+r)=Ar$ for a linear map
$A:\R^2\to q_o^\perp$, then
\[
 M_{r,s}=\frac{r\cdot s}{2}\sum_{i=1}^2|Ae_i|^2,
 \qquad
 S_o=\frac12\sum_{i=1}^2|Ae_i|^2,
\]
so $Q_{r,s}$ vanishes on every affine tangent field.  This uses the
symmetric part of the quarter-turn average; full dihedral symmetry is
not required.

A quadratic label is the entire compensated packet $aQ_{r,s}$, with
its compensating nearest-neighbor coefficients tagged to that packet.
Its norm contribution is
\begin{equation}
 |a|\left\{W_\sigma(1+|r|_1+|s|_1)
                  +|r\cdot s|W_\sigma(3)\right\}.
 \label{eq:setup-quadratic-norm}
\end{equation}
The norm of a quadratic list is the supremum over anchors of the sum of
these contributions.  These norms are norms on the retained tagged
coefficient lists, not on the resulting polynomial after all
cancellations: different retained lists can represent the same
polynomial.  Thus a compensating coefficient uses its own
nearest-neighbor path weight, while its tag retains the original
packet's cancellation and support information.  This convention never
shortens a mask or a complete record.  Equations
\eqref{eq:setup-canonical-norm} and \eqref{eq:setup-quadratic-norm} are the
fixed norms of \Rref{rg:canonical-norm} and
\Rref{supp:quadratic-packet-norm}, with $\mathcal C_4$ replacing the
dihedral group.

The retained coefficient bounds are
\begin{equation}
 \|P_4\|_\sigma\le B_1,\quad
 \|I_2\|_\sigma\le B_2,\quad
 \|P_6\|_\sigma\le B_5,\quad
 \|I_4\|_\sigma\le B_6,\quad
 \|J_2\|_\sigma\le B_7.
 \label{eq:setup-canonical-caps}
\end{equation}
The caps are fixed successively in the indicated slot order.  Polynomial
formulas and their affine compensation are formed in the bulk before
folding onto a torus.  A label with $u\le\mathfrak m_j$ may use a common
mask ball of radius $C_\chi\mathfrak m_j$ about its anchor.  Longer paths
are padded by fixed multiples of $\mathfrak m_j$.  The resulting load is
at most $C(1+u/\mathfrak m_j)$, with the full compensated packet retaining
the required graph.

\subsubsection{Differentiable errors}

The coefficient norms control the polynomial part of the density.  The
remaining regular functions are controlled on the graph domains
\eqref{eq:setup-graph-domain}, including enough derivatives to compare
two successive exact integrations.

A regular-error label at $o$ consists of a function
$f_X\in C^8(\mathcal U_{X,j})$, a graph $X$, and its complete record of
load $s_X$.  Write $v\times$ for the skew-symmetric linear map
$w\mapsto v\times w$ on $\R^3$.  For $1\le r\le8$, define the rotation
jet
\begin{equation}
 D_{v_1,\ldots,v_r}f_X(q)
 =\left.
  \partial_{\theta_1}\cdots\partial_{\theta_r}
  f_X\left(\left(
    \exp\!\left[\left(\sum_{a=1}^r\theta_av_{a,x}\right)\times\right]
                     q_x\right)_x\right)
 \right|_{\theta=0}.
 \label{eq:setup-error-jets}
\end{equation}
In each derivative slot independently, take all vector fields satisfying
\begin{equation}
 |v_{a,x}|\le t_j(1+\operatorname{dist}(x,o))^8.
 \label{eq:setup-error-directions}
\end{equation}
Distances in this bound are physical torus distances on a finite torus.
For $0\le k\le8$, let
\begin{equation}
 [f_X]_{k,j}
 =\sum_{r=0}^k
   \sup_{q\in\mathcal U_{X,j}}
   \sup_{v_1,\ldots,v_r}
       |D_{v_1,\ldots,v_r}f_X(q)|,
 \label{eq:setup-error-seminorm}
\end{equation}
where the $r=0$ term is $\sup_{\mathcal U_{X,j}}|f_X|$.
These are classical derivatives on the open graph domain.  Values away
from that domain require no differentiability.  This makes explicit the
rotation version of \Rref{rg:directions}.

Set $A=64$.  The error-list norm and its cap are
\begin{equation}
 \|f\|_{k,j,A}
   =\sup_o\sum_{X\text{ anchored at }o}e^{As_X}[f_X]_{k,j},
 \qquad
 \|f\|_{8,j,A}\le\delta_j:=H_j^{-2.05}.
 \label{eq:setup-error-list-norm}
\end{equation}
Each regular function is invariant under simultaneous $O(3)$
transformations of its spin arguments.  Each bulk packet is also
averaged under the quarter turns about its anchor, including position
inversion $x-o\mapsto-(x-o)$, after its masks have been strengthened to
a common invariant graph.  Short folded copies inherit this inversion
parity from their lifted bulk packets.

For a smooth local scalar $F$ with these spin and spatial symmetries,
its \emph{affine coefficient} at $o$ is
\begin{equation}
 \mathfrak h_o(F)=
 \left.\frac{d^2}{d\theta^2}
 F\left(\left(\operatorname{Exp}_n
       \bigl(\theta v\,\widehat e\cdot(x-o)\bigr)\right)_x\right)
 \right|_{\theta=0},
 \label{eq:setup-affine-coefficient}
\end{equation}
where $n\in S^2$, $v\in n^\perp$ is a unit vector, and $\widehat e$ is
a coordinate unit vector in the current lattice frame.  Symmetry makes
this number independent of these choices.  It is an ordinary second
derivative, so $\mathfrak h_o(S_o)=1$.  Polarization gives the mixed
affine Hessian
$\mathfrak h_o(F)\operatorname{Tr}(A_1^*A_2)$ for tangent maps
$A_1,A_2:\R^2\to n^\perp$.  The affine coefficient of a bulk list means
the sum of these numbers per anchor.

A bulk or short error has zero value at every aligned configuration and
zero affine Hessian there.  More precisely, fix $n\in S^2$ and linear
maps $A_1,A_2:\R^2\to n^\perp$.  With the anchor kept at $n$, its
normalizations are
\begin{equation}
 f_X((n)_x)=0,
 \qquad
 \left.\partial_s\partial_t
 f_X\left(\left(
   \operatorname{Exp}_{n}\bigl(sA_1(x-o)+tA_2(x-o)\bigr)
                            \right)_x\right)
 \right|_{s=t=0}=0.
 \label{eq:setup-error-normalization}
\end{equation}
These affine tests are made on the lifted bulk formula and do not have
to be periodic.  The first derivative at alignment vanishes by spin
invariance.  Winding errors are allowed to have nonzero values and
affine Hessians at alignment.  These are the normalization conventions
of \Rref{rg:error-norm} and \Rref{rg:normalization}.

\subsubsection{The mandatory covering sum}

A covering label $\ell$ consists of a bounded measurable function
$k_\ell(q)$, a complete support $P_\ell$, a record of load $s_\ell$, and
a nonempty inventory $J_\ell$ of edges whose endpoints lie in that
support.  It is required that
\[
 k_\ell(q)=0\quad\text{unless}\quad J_\ell\subset D(q).
\]
Thus $J_\ell$ is the inventory denoted by $D_\ell$ in
\Rref{rg:covering-gas}.
No sign or differentiability condition is imposed on $k_\ell$.
Define
\begin{equation}
 \Xi(q)=
  \sum_{\substack{\Gamma\text{ a finite family of labels}\\
                   P_\ell\cap P_{\ell'}=\varnothing\ (\ell\ne\ell')\\
                   \bigsqcup_{\ell\in\Gamma}J_\ell=D(q)}}
                       \prod_{\ell\in\Gamma}k_\ell(q).
 \label{eq:setup-covering-sum}
\end{equation}
The empty family has weight one.  In particular $\Xi(q)=1$ when
$D(q)=\varnothing$.  If $D(q)$ is nonempty, every one of its edges must
be supplied exactly once by the inventories in the family; inventories
are not optional decorations.  The covering norm is
\begin{equation}
 \|k\|_{j,A}
   =\sup_x\sum_{\ell:x\in P_\ell}
                e^{As_\ell}\|k_\ell\|_\infty,
 \qquad
 \|k\|_{j,A}\le w_j:=\exp(-p_j^{1/4}).
 \label{eq:setup-covering-norm}
\end{equation}
The supremum norm is over the spin arguments of the weight, with all
of its recorded tests included.  This is \Rref{rg:covering-gas}.
The covering sum is absolutely convergent even for a countable label
list.  Indeed, dropping compatibility and inventory constraints and
summing over all finite subsets of the label list gives the bound
\begin{equation}
 \sum_{\Gamma}\prod_{\ell\in\Gamma}\|k_\ell\|_\infty
 \le\exp\!\left(\sum_\ell\|k_\ell\|_\infty\right)
 \le\exp(|\Lambda_j|w_j).
 \label{eq:setup-covering-convergence}
\end{equation}
For the second inequality, each support is nonempty, so
\[
 \sum_\ell\|k_\ell\|_\infty
 \le\sum_{x\in\Lambda_j}\sum_{\ell:x\in P_\ell}\|k_\ell\|_\infty
 \le|\Lambda_j|w_j.
\]
In particular $|\Xi(q)|\le\exp(|\Lambda_j|w_j)$.

All lists are covariant under simultaneous $O(3)$ transformations of
the spin arguments, with constituent choices and tags retained together.
The individual regular functions have the invariance already required
above.  Anchored lists are translation covariant in the bulk.  Their
spatial rules commute with quarter turns, and with simultaneous
reflection of the blocking geometry and its data.  In particular a
reflection can exchange the two inclined blocking types.  On an
asymmetric torus, the symmetry operations concern lifted bulk labels
and their folded copies; they do not require the torus itself to have a
quarter-turn symmetry.  Any operation encountering a long record or a
period disagreement carries the winding load.  Tags and full invariant
tensors are retained together under these operations.

Anchors are assigned equivariantly by distributing a term equally among
its permitted anchors while retaining identical complete supports.
Eligibility and internal compatibility tests remain part of their
labels.  When masks in a spatial orbit differ, they are first replaced
by a common stronger graph mask.  The difference is retained exactly
as a covering correction, since it can be nonzero only when that graph
contains an edge in $D(q)$.  The same convention applies to all later
mask strengthening; its algebraic reassembly is the one in
\Rref{rg:connected}.

With these definitions, an effective density relative to product
normalized area measure on $(S^2)^{\Lambda_j}$ has the form
\begin{equation}
 \begin{aligned}
 \rho_j(q)={}&\exp\{\mathfrak v|\Lambda_j|-b\mathcal E_{h,t_j}(q)\}\Xi(q)\\
 &\quad\times
 \exp\left\{-\sum_{\text{canonical labels }\alpha}
                         1_{X_\alpha}(q)\,\mathcal P_\alpha(q)
       -\sum_X1_X(q)f_X(q)\right\}.
 \end{aligned}
 \label{eq:setup-effective-density}
\end{equation}
Here $\mathcal P_\alpha$ includes the power of $b$ specified by its
slot in \eqref{eq:setup-slots}.  The bulk scalar $\mathfrak v$ is the same throughout the compatible
family of representations over admitted periods.  Constants caused
by a period-dependent discrepancy remain in winding errors.  This is
the $S^2$ version of \Rref{rg:density}.  Estimates for the representation
also permit signed densities; densities obtained by the actual
observation kernel are positive.  We call a representation \emph{admitted}
when these support, symmetry, period, and norm conditions hold and
$b\in[H_j/2,2H_j]$.

\subsubsection{Comparing two densities}

Two inputs are compared on a common reference layer, with the same
$H_j,t_j$ and geometric thresholds.  First use the prescribed common
refinements of their labels, masks, and complete records; any mask
strengthening retains its exact covering correction as above.  Then
retain the union of these refined label lists and assign coefficient or
weight zero when a label is absent.  Corresponding functions now have
common graph domains and complete records.  Write $b_1,b_2$ for the precise
couplings and set $\lambda=b_2-b_1$.  For fixed positive weights
$\omega_1,\ldots,\omega_7,\omega_f,\omega_k$, define
\begin{equation}
 u=\sum_{a=1}^7\omega_a
          \|\mathbf P_{2,a}-\mathbf P_{1,a}\|_\sigma
    +\omega_f\delta_j^{-1}\|f_2-f_1\|_{7,j,A}
    +\omega_kw_j^{-1}\|k_2-k_1\|_{j,A}.
 \label{eq:setup-density-discrepancy}
\end{equation}
The third and fourth coefficient discrepancies are zero.  Polynomial
coefficients are compared without their powers of $b$; the change in
the precise coupling is recorded by $\lambda$.  The volume scalar is
tracked separately and is not a component of $u$.

The weights are chosen after the successive canonical caps so that the
triangular coefficient estimate and the error and covering estimates
hold together.  They remain fixed across layers and admitted periods.
Each individual error is controlled through eight derivatives, whereas
an error discrepancy is measured through seven.  Thus a difference of
integration maps acting on an unchanged error can use its eighth
derivative without reducing the regularity of the next individual
density.  This is the comparison convention of \Rref{rg:closure}, with
the stronger contraction proved in Section~\ref{sec:rg}.


\subsection{Endpoint laws and the scope of the imported estimates}
\label{setup:endpoints}

An \emph{endpoint density} is a layer-zero representation satisfying the
preceding caps, with \(b\in[H/2,2H]\), and with the bulk scalar removed.
We denote its remaining exponent by \(X\).  For an actual observed
representation, \(e^X\Xi>0\); its partition function and probability
law are therefore well defined by
\begin{equation}
 \begin{split}
 X(q)&=-b\mathcal E_{h,t_0}(q)
       -\sum_\alpha1_{X_\alpha}(q)\mathcal P_\alpha(q)
       -\sum_Y1_Y(q)f_Y(q),\\
 Z&=\int e^{X(q)}\Xi(q)\,\dd\omega^{\otimes\Lambda_0}(q),
 \qquad
 \dd\mu(q)=Z^{-1}e^{X(q)}\Xi(q)\,\dd\omega^{\otimes\Lambda_0}(q).
 \end{split}
 \label{setup:endpoint-law}
\end{equation}
The scalar removed from the
partition function is exactly \(e^{\mathfrak v|\Lambda_0|}\).
For a fixed endpoint reference scale, \(\delta_0=H^{-2.05}\) and
\(w_0=\exp[-(\log H)^{P_0/4}]\); with \(P_0>4\), the latter tends
to zero faster than every fixed inverse power of \(H\).

The definitions distinguish three kinds of input to the exact blocking
construction.
The scalar identities and localization estimates of
\Rref{free:bounds} and \Rref{lem:free-stack} have no spin-dimension
hypothesis.  Their extension to the extra histories is proved in
Section~\ref{sec:free}.  The Gaussian product and connected-expansion
estimates of \Rref{supp:gaussian-product} and \Rref{rg:tree-condition}
apply to finite-dimensional Gaussian integrations and complete support
records satisfying their stated envelope and support-hit bounds.
Section~\ref{sec:rg} verifies those bounds for the actual \(S^2\)
factors, including derivatives and the auxiliary normal coordinate.
Finally, the improved coefficient, error, endpoint, and source
comparisons are conclusions of this paper.  They are not included in
the definition of an endpoint law.  The separate adaptations of the
preliminary mixing estimates and the finite-volume Laplace coefficients
are stated at their points of use in Sections~\ref{sec:preliminary}
and~\ref{sec:shooting}.

This separation is useful in two places.  Signed covering weights can
be estimated algebraically without assuming positivity of a truncated
covering sum.  Reflection positivity and the continuum spectral
conclusion, by contrast, concern the complete positive laws and are
established only after the corresponding observation and limit
arguments.
